\section{Research Roadmap}
\label{sec:roadmap}

Table~\ref{tab:roadmap} presents a research roadmap organized around the ICA six-layer model, spanning three phases: short-term (1--2 years), mid-term (2--4 years), and long-term (4--8 years). Each entry specifies the problem to be solved, the proposed method, and the expected outcome, making every item a concrete, actionable research project. Figure~\ref{fig:roadmap_swimlane} provides a visual overview of the roadmap as a swim-lane diagram organized by ICA layer and time horizon.

\footnotesize
\renewcommand{\arraystretch}{1.35}
\begin{longtable}{>{\small}p{0.06\textwidth}>{\small}p{0.06\textwidth}>{\small}p{0.13\textwidth}>{\small}p{0.25\textwidth}>{\small}p{0.18\textwidth}>{\small}p{0.14\textwidth}}
\caption{Research agenda organized by ICA layer, phase, and associated law or axiom.}
\label{tab:roadmap}\\
\toprule
\textbf{Phase} & \textbf{ICA Layer} & \textbf{Topic} & \textbf{Method} & \textbf{Metrics} & \textbf{Related Law / Axiom} \\
\midrule
\endfirsthead
\toprule
\textbf{Phase} & \textbf{ICA Layer} & \textbf{Topic} & \textbf{Method} & \textbf{Metrics} & \textbf{Related Law / Axiom} \\
\midrule
\endhead

Short-term & L2 & KV cache hit-rate optimization & Semantic eviction; KV quantization; prefix sharing & $H$, $S$, TTFT & Law~I (Semantic Locality) \\[3pt]
Short-term & L3 & Unified context compiler & Hot / warm / cold tiering; version stamps with TTL & $W_{\mathrm{eff}}$, compression ratio & Law~II (Context Budget) \\[3pt]
Short-term & L4 & Tool ABI standardization & Schema versioning; capability annotation & Invocation success rate; security violation rate & Axiom~5 (Least Privilege) \\[3pt]
Short-term & L5 & Software engineering agent kernel & Unified ACI; specialized sub-agents & Task completion rate; $E$ & Law~III (Agent Speedup) \\[3pt]
Mid-term & L3 & Semantic page replacement & Learned eviction; summary write-back retrieval & $\beta(L)$ curve & Law~II \\[3pt]
Mid-term & L5 & Consistency-controlled shared memory & Event sourcing; conflict detection & Memory correctness rate & Axiom~4 (Virtualization) \\[3pt]
Mid-term & L2 & Heterogeneous model collaboration & Speculative decoding; MoE routing & End-to-end latency; $H$ & Law~I \\[3pt]
Mid-term & L4 & Agent security architecture & Sandboxing; audit logging; capability security & Audit completeness & Axioms~5,~6 \\[3pt]
Long-term & L1--L6 & Intelligent ISA and programmable skill layer & Task DSL; controlled output formats; reusable skill packages & Generalization ability & Axioms~2,~3 \\[3pt]
Long-term & L3--L5 & Stateful individual agents & Long-term memory; experience abstraction and strategy learning & Experience utilization rate & Laws~II,~III \\[3pt]
Long-term & L1--L6 & Distributed model-native compute fabric & Memory replication; consistency layers; load balancing & Cluster utilization & Laws~I,~III \\[3pt]
Long-term & L5 & Governance-layer primitives & Probabilistic / deterministic dual-plane; audit logs; permission declarations & Governance coverage & Axioms~3,~5,~6 \\[3pt]
Long-term & L1--L2 & Substrate-independent compute stack & Non-silicon interface generalization; baseline intelligence score threshold & Interface compatibility rate & Axiom~2 \\[3pt]
Long-term & L4--L5 & Physical-world interface drivers & Sensor / actuator driver abstraction; simulation environment integration & Driver coverage rate & Axioms~2,~5 \\[3pt]
Mid-term & Full stack & Weakest-link evaluation framework & Worst-at-$k$; multi-dimensional capability profiling; compliance test suite & Minimum dimension score & Barrel Principle \\[3pt]
\bottomrule
\end{longtable}
\renewcommand{\arraystretch}{1.0}
\normalsize

% ---------------------------------------------------------------------------
% Research Roadmap Swim-Lane Diagram  (Gantt-style)
% fig:roadmap_swimlane
% ---------------------------------------------------------------------------
\begin{figure}[htbp]
\centering
% ---- color definitions (outside tikzpicture so caption can use them) ----
\providecolor{lawI}{HTML}{4C78A8}       % blue   – inference efficiency
\providecolor{lawII}{HTML}{D9824B}      % orange – context management
\providecolor{lawIII}{HTML}{C44E52}     % red    – agent speedup
\providecolor{axiom}{HTML}{8C8C8C}      % gray   – architectural principles
%
\resizebox{0.97\textwidth}{!}{%
\begin{tikzpicture}[
    % ---- global styles ----
    >=Latex,
    font=\small,
    lane label/.style={
        font=\footnotesize\bfseries, anchor=east, text=black!75
    },
    phase header/.style={
        font=\small\bfseries, anchor=south, text=black!80
    },
    phase sub/.style={
        font=\scriptsize, anchor=north, text=black!50
    },
    bar/.style={
        draw=#1!60!black, fill=#1!25, rounded corners=3pt,
        minimum height=0.55cm, inner xsep=3pt, inner ysep=1.5pt,
        font=\scriptsize, align=center, line width=0.5pt,
        text=black!85
    },
    bar selected/.style={
        draw=#1!80!black, fill=#1!35, rounded corners=3pt,
        minimum height=0.55cm, inner xsep=3pt, inner ysep=1.5pt,
        font=\scriptsize\bfseries, align=center, line width=0.7pt,
        text=black!90
    },
    dep arrow/.style={
        -{Latex[length=1.8mm]}, thick, black!45, line width=0.6pt
    },
    dep arrow dashed/.style={
        -{Latex[length=1.8mm]}, thick, black!30,
        line width=0.5pt, dashed
    },
]

% ================================================================
% COLOR PALETTE  (lane backgrounds only; main colors defined above)
% ================================================================
\providecolor{laneBg}{HTML}{F8F9FA}     % very light gray for lane fills
\providecolor{laneAlt}{HTML}{F0F2F5}    % slightly darker for alternating

% ================================================================
% LAYOUT PARAMETERS
% ================================================================
% x-coordinates for three phase columns
\def\xShortL{3.2}
\def\xShortR{7.2}
\def\xMidL{8.2}
\def\xMidR{12.8}
\def\xLongL{13.8}
\def\xLongR{19.4}

% y-coordinates for six lanes (L1 at bottom, L6 at top)
\def\yLone{0.5}
\def\yLtwo{2.0}
\def\yLthree{3.5}
\def\yLfour{5.0}
\def\yLfive{6.5}
\def\yLsix{8.0}

% lane vertical extent
\def\laneHalf{0.62}

% ================================================================
% PHASE HEADER LABELS
% ================================================================
\node[phase header] at ({(\xShortL+\xShortR)/2}, \yLsix+1.45)
    {Short-Term};
\node[phase sub] at ({(\xShortL+\xShortR)/2}, \yLsix+1.05)
    {1--2 years};

\node[phase header] at ({(\xMidL+\xMidR)/2}, \yLsix+1.45)
    {Mid-Term};
\node[phase sub] at ({(\xMidL+\xMidR)/2}, \yLsix+1.05)
    {2--4 years};

\node[phase header] at ({(\xLongL+\xLongR)/2}, \yLsix+1.45)
    {Long-Term};
\node[phase sub] at ({(\xLongL+\xLongR)/2}, \yLsix+1.05)
    {4--8 years};

% ================================================================
% PHASE BACKGROUND BANDS (alternating subtle shading)
% ================================================================
\fill[laneBg, rounded corners=4pt]
    (\xShortL-0.2, \yLone-\laneHalf-0.15)
    rectangle (\xShortR+0.2, \yLsix+\laneHalf+0.15);
\fill[laneAlt, rounded corners=4pt]
    (\xMidL-0.2, \yLone-\laneHalf-0.15)
    rectangle (\xMidR+0.2, \yLsix+\laneHalf+0.15);
\fill[laneBg, rounded corners=4pt]
    (\xLongL-0.2, \yLone-\laneHalf-0.15)
    rectangle (\xLongR+0.2, \yLsix+\laneHalf+0.15);

% Phase separator vertical lines
\draw[black!15, line width=0.4pt]
    ({(\xShortR+\xMidL)/2}, \yLone-\laneHalf-0.15)
    -- ({(\xShortR+\xMidL)/2}, \yLsix+\laneHalf+0.15);
\draw[black!15, line width=0.4pt]
    ({(\xMidR+\xLongL)/2}, \yLone-\laneHalf-0.15)
    -- ({(\xMidR+\xLongL)/2}, \yLsix+\laneHalf+0.15);

% ================================================================
% LANE HORIZONTAL GUIDES  (dashed, very subtle)
% ================================================================
\foreach \y in {\yLone,\yLtwo,\yLthree,\yLfour,\yLfive,\yLsix}{
    \draw[black!8, line width=0.3pt, dashed]
        (\xShortL-0.2, \y) -- (\xLongR+0.2, \y);
}

% ================================================================
% LANE LABELS  (L1–L6, left margin)
% ================================================================
\node[lane label] at (\xShortL-0.55, \yLsix) {L6 Application};
\node[lane label] at (\xShortL-0.55, \yLfive) {L5 Orchestration};
\node[lane label] at (\xShortL-0.55, \yLfour) {L4 Interface};
\node[lane label] at (\xShortL-0.55, \yLthree) {L3 Context};
\node[lane label] at (\xShortL-0.55, \yLtwo) {L2 Inference};
\node[lane label] at (\xShortL-0.55, \yLone) {L1 Physical};

% ================================================================
% "WE ARE HERE" VERTICAL LINE
% ================================================================
\draw[black!70, line width=1.1pt, densely dashed]
    (\xShortL-0.05, \yLone-\laneHalf-0.15)
    -- (\xShortL-0.05, \yLsix+\laneHalf+0.15);
\pgfmathsetmacro{\weAreHereY}{(\yLone+\yLsix)/2}
\node[font=\scriptsize\bfseries, text=black!65, rotate=90, anchor=south]
    at (\xShortL-0.45, \weAreHereY)
    {We Are Here};

% ================================================================
% RESEARCH TOPIC BARS
% ================================================================

% ---- L6  (Application) --------------------------------------------------
\node[bar=axiom]
    (L6gov) at ({(\xMidL+\xMidR)/2}, \yLsix)
    {Governance kernel};

\node[bar=lawIII]
    (L6dist) at ({(\xLongL+\xLongR)/2-0.2}, \yLsix)
    {Distributed agent};
\node[bar=axiom]
    (L6fed) at ({(\xLongL+\xLongR)/2-0.2}, \yLsix+0.52)
    {Federated agent networks};

% ---- L5  (Orchestration) ------------------------------------------------
\node[bar selected=lawIII]
    (L5proto) at ({(\xShortL+\xShortR)/2}, \yLfive)
    {Protocol unification};

\node[bar=lawIII]
    (L5mcp) at ({(\xMidL+\xMidR)/2-0.55}, \yLfive-0.28)
    {MCP + A2A integration};
\node[bar=axiom]
    (L5cap) at ({(\xMidL+\xMidR)/2+0.55}, \yLfive+0.28)
    {Capability negotiation};

\node[bar=lawIII]
    (L5os) at ({(\xLongL+\xLongR)/2}, \yLfive)
    {Agent OS kernel};

% ---- L4  (Semantic Interface) -------------------------------------------
\node[bar=lawIII]
    (L4bench) at ({(\xShortL+\xShortR)/2}, \yLfour)
    {Agent runtime\\[-1pt]benchmarks};

\node[bar=lawI]
    (L4interop) at ({(\xMidL+\xMidR)/2-0.55}, \yLfour-0.28)
    {Cross-framework\\[-1pt]interop};
\node[bar=axiom]
    (L4sec) at ({(\xMidL+\xMidR)/2+0.55}, \yLfour+0.28)
    {Agent security\\[-1pt]architecture};

\node[bar=axiom]
    (L4driver) at ({(\xLongL+\xLongR)/2-0.3}, \yLfour-0.28)
    {Physical-world\\[-1pt]interface drivers};
\node[bar=axiom]
    (L4isa) at ({(\xLongL+\xLongR)/2-0.3}, \yLfour+0.28)
    {Intelligent ISA};

% ---- L3  (Context Management) -------------------------------------------
\node[bar selected=lawII]
    (L3comp) at ({(\xShortL+\xShortR)/2}, \yLthree)
    {Context compilation\\[-1pt]algorithms};

\node[bar=lawII]
    (L3tier) at ({(\xMidL+\xMidR)/2-0.55}, \yLthree-0.28)
    {Semantic page\\[-1pt]replacement};
\node[bar=axiom]
    (L3mem) at ({(\xMidL+\xMidR)/2+0.55}, \yLthree+0.28)
    {Consistency-controlled\\[-1pt]shared memory};

\node[bar=lawII]
    (L3persist) at ({(\xLongL+\xLongR)/2}, \yLthree)
    {Persistent context\\[-1pt]stores (L3--L5)};

% ---- L2  (Inference Serving) --------------------------------------------
\node[bar selected=lawI]
    (L2kv) at ({(\xShortL+\xShortR)/2-0.2}, \yLtwo-0.28)
    {KV cache\\[-1pt]optimization};
\node[bar=lawI]
    (L2sparse) at ({(\xShortL+\xShortR)/2-0.2}, \yLtwo+0.28)
    {Sparse attention\\[-1pt]optimization};

\node[bar=lawI]
    (L2prefix) at ({(\xMidL+\xMidR)/2-0.55}, \yLtwo-0.28)
    {Heterogeneous model\\[-1pt]collaboration};
\node[bar=lawI]
    (L2semcache) at ({(\xMidL+\xMidR)/2+0.55}, \yLtwo+0.28)
    {Prefix caching\\[-1pt]standardization};

\node[bar=axiom]
    (L2substrate) at ({(\xLongL+\xLongR)/2}, \yLtwo)
    {Substrate-independent\\[-1pt]compute stack};

% ---- L1  (Physical Execution) -------------------------------------------
\node[bar=lawI]
    (L1sparse) at ({(\xShortL+\xShortR)/2}, \yLone)
    {Sparse attention\\[-1pt]kernels};

\node[bar=axiom]
    (L1alt) at ({(\xLongL+\xLongR)/2}, \yLone)
    {Alternative inference\\[-1pt]substrates};

% ================================================================
% DEPENDENCY ARROWS
% ================================================================

% KV cache optimization -> Context compilation (L2 -> L3)
\draw[dep arrow] (L2kv.north) -- (L3comp.south);

% Context compilation -> Semantic page replacement (L3 short -> L3 mid)
\draw[dep arrow] (L3comp) -- (L3tier);

% Protocol unification -> MCP+A2A integration (L5 short -> L5 mid)
\draw[dep arrow] (L5proto) -- (L5mcp);

% MCP+A2A integration -> Capability negotiation (L5 mid -> L5 mid)
\draw[dep arrow] (L5mcp) -- (L5cap);

% Agent runtime benchmarks -> Cross-framework interop (L4 short -> L4 mid)
\draw[dep arrow] (L4bench) -- (L4interop);

% Agent security architecture -> Governance kernel (L4 mid -> L6 mid)
\draw[dep arrow] (L4sec.north) -- (L6gov.south);

% Governance kernel -> Distributed agent coordination (L6 mid -> L6 long)
\draw[dep arrow] (L6gov) -- (L6dist);

% Consistency shared memory -> Persistent context stores (L3 mid -> L3 long)
\draw[dep arrow] (L3mem) -- (L3persist);

% Capability negotiation -> Agent OS kernel (L5 mid -> L5 long)
\draw[dep arrow dashed] (L5cap.east) -- (L5os.west);

% Intelligent ISA -> Agent OS kernel (L4 long -> L5 long)
\draw[dep arrow] (L4isa.north) -- (L5os.south);

% Substrate-independent stack -> Alternative inference substrates (L2 long -> L1 long)
\draw[dep arrow] (L2substrate.south) -- (L1alt.north);

% Cross-framework interop -> Agent security architecture (L4 mid -> L4 mid)
\draw[dep arrow] (L4interop.north east) -- (L4sec.south west);

% Semantic page replacement -> Persistent context stores (L3 mid -> L3 long)
\draw[dep arrow dashed] (L3tier.east) -- (L3persist.west);

% Heterogeneous model collaboration -> Prefix caching (L2 mid -> L2 mid)
\draw[dep arrow] (L2prefix.north east) -- (L2semcache.south west);

% Distributed agent -> Federated agent networks (L6 long -> L6 long)
\draw[dep arrow] (L6dist) -- (L6fed);

% Physical-world interface drivers -> Agent OS kernel (L4 long -> L5 long)
\draw[dep arrow dashed] (L4driver.north) -- (L5os.south west);

% ================================================================
% COLOR LEGEND
% ================================================================
\pgfmathsetmacro{\legY}{-1.05}
\pgfmathsetmacro{\xA}{(\xShortL+\xShortR)/2}
\pgfmathsetmacro{\xB}{(\xMidL+\xMidR)/2}
\pgfmathsetmacro{\xC}{(\xLongL+\xLongR)/2}

\node[font=\scriptsize\bfseries, anchor=east, text=black!65]
    at (\xA-0.3, \legY) {Legend:};

% Law I
\fill[lawI!25, draw=lawI!60!black, rounded corners=2pt, line width=0.5pt]
    (\xA, \legY-0.17) rectangle ++(1.6, 0.34);
\node[font=\scriptsize, anchor=west, text=black!70]
    at (\xA+1.8, \legY) {Law~I (Semantic Locality)};

% Law II
\fill[lawII!25, draw=lawII!60!black, rounded corners=2pt, line width=0.5pt]
    (\xB-0.5, \legY-0.17) rectangle ++(1.6, 0.34);
\node[font=\scriptsize, anchor=west, text=black!70]
    at (\xB+1.3, \legY) {Law~II (Context Budget)};

% Law III
\fill[lawIII!25, draw=lawIII!60!black, rounded corners=2pt, line width=0.5pt]
    (\xB+3.2, \legY-0.17) rectangle ++(1.6, 0.34);
\node[font=\scriptsize, anchor=west, text=black!70]
    at (\xB+5.0, \legY) {Law~III (Agent Speedup)};

% Axiom
\fill[axiom!25, draw=axiom!60!black, rounded corners=2pt, line width=0.5pt]
    (\xC-1.0, \legY-0.17) rectangle ++(1.6, 0.34);
\node[font=\scriptsize, anchor=west, text=black!70]
    at (\xC+0.8, \legY) {Axiom (Architectural)};

\end{tikzpicture}}%
\caption{Research roadmap organized by ICA layer and time horizon.
  Horizontal lanes correspond to ICA layers L1--L6; columns represent
  short-term (1--2\,yr), mid-term (2--4\,yr), and long-term (4--8\,yr)
  research phases.  Bar colors indicate which law or axiom each topic
  validates: {\color{lawI}Law~I} (Semantic Locality, blue),
  {\color{lawII}Law~II} (Context Budget, orange),
  {\color{lawIII}Law~III} (Agent Speedup, red), and
  {\color{axiom}Axiom} (architectural principles, gray).
  Arrows denote key dependencies between research topics.}
\label{fig:roadmap_swimlane}
\end{figure}


\subsection{Short-Term Research Topics}

The short-term items focus on optimizing and standardizing existing engineering prototypes, with the goal of bringing current systems to production-grade performance and reliability.

\paragraph{KV cache hit-rate optimization (L2).}
The decoding phase of LLM inference is a classic memory-bandwidth-bound problem~\cite{gholami2024memorywall}, and the KV cache hit rate $H$ directly governs the inference speedup ratio $S$ predicted by the Semantic Locality Law (Law~I).
Three complementary techniques can be brought to bear.
First, \textit{semantic eviction} strategies, exemplified by H2O's heavy-hitter identification~\cite{h2o2023}, retain the KV entries most likely to be reused based on attention weight patterns rather than recency alone.
Second, \textit{KV quantization} compresses the KV cache with minimal quality loss: KVQuant achieves up to $8\times$ compression~\cite{hooper2024kvquant}, while KIVI introduces asymmetric 2-bit quantization~\cite{liu2024kivi}.
Third, \textit{prefix sharing} across requests, as implemented by PagedAttention~\cite{kwon2023pagedattention} and RadixAttention~\cite{sglang2024}, enables multiple requests with common system prompts to reuse the same KV blocks, directly increasing $H$.
The target outcome is to raise $H$ from the current 0.5--0.7 range to above 0.85 under typical agent workloads, yielding the 5--10$\times$ speedup predicted by Law~I.

\paragraph{Unified context compiler (L3).}
Different agent frameworks manage context in idiosyncratic ways: MemGPT uses virtual context management~\cite{memgpt2023}, Letta builds stateful agents with explicit memory tiers~\cite{letta_stateful_2026}, and HiAgent employs hierarchical working memory~\cite{hiagent2024}.
What is missing is a unified \textit{context compiler}, analogous to an operating system's page replacement subsystem, that decides, for each piece of information, whether to load it verbatim, as a summary, or as a retrievable index entry.
We propose a hot / warm / cold tiered context management strategy, where each context block carries a version stamp and a time-to-live (TTL) annotation.
The expected result is a $\geq 50\%$ increase in the effective context utilization ratio $W_{\mathrm{eff}}/C$ without degrading task completion rates, thereby providing direct empirical validation of the Context Budget Law (Law~II).

\paragraph{Tool ABI standardization (L4).}
While MCP and A2A have made significant strides at the protocol level~\cite{mcp_intro_2026,agentprotocols2025}, tool schemas themselves still lack unified specifications for versioning, backward compatibility, and capability annotation.
We propose enriching tool schemas with semantic version numbers, input/output type constraints, side-effect declarations, and capability annotations, transforming L4 into a genuine Application Binary Interface (ABI) for the model-native computing stack.
The expected outcome is a measurable increase in tool invocation success rates and a decrease in security violations, providing the interface foundation required to operationalize Axiom~5 (Least Privilege).

\paragraph{Software engineering agent kernel (L5).}
SWE-bench and SWE-agent have demonstrated that current code agents still achieve limited resolution rates on real-world GitHub issues~\cite{swebench2023,sweagent2024}.
A key bottleneck is the lack of a unified Agent-Computer Interface (ACI) standard: different frameworks adopt inconsistent sub-agent delegation strategies.
We propose standardizing the ACI specification and designing specialized sub-agents (a code-search agent, an edit agent, and a test agent) whose orchestration efficiency $E$ can be quantified through the Agent Speedup Law (Law~III).
The target is to increase task completion rates on SWE-bench-class benchmarks while raising $E$ from the current 0.3--0.5 range to above 0.6.

\subsection{Mid-Term Research Topics}

The mid-term items shift focus from per-layer optimization to cross-layer coordination and system-level security, moving model-native computing from ``point optimizations'' to ``full-stack co-design.''

\paragraph{Semantic page replacement (L3).}
Classical operating systems replace pages using deterministic access patterns (LRU, CLOCK). In a model-native computing system, ``page replacement'' must be governed by \textit{semantic relevance}, deciding which context blocks to evict based on their pertinence to the current task.
We propose a learned eviction policy that scores each context block against the active task and generates compressed summaries for evicted blocks, enabling later retrieval.
The primary research output is an empirically measured $\beta(L)$ curve that characterizes how attention retention decays with position, providing the Context Budget Law with calibrated quantitative parameters.

\paragraph{Consistency-controlled shared memory (L5).}
When multiple agents collaborate, shared memory consistency management lacks a unified framework.
Concurrent modifications to the same memory block can cause conflicts and information loss, the analog of cache coherence problems in multiprocessor systems.
We propose introducing event sourcing and conflict detection mechanisms, combined with tunable consistency levels ranging from eventual consistency to strong consistency~\cite{raft2014,cap2002}.
The target is to ensure that multi-agent memory correctness rates match single-agent baselines, providing empirical grounding for Axiom~4 (Virtualization) at the memory layer.

\paragraph{Heterogeneous model collaboration (L2).}
Models of different sizes offer distinct trade-offs among latency, cost, and capability, yet current systems typically deploy a single model.
Heterogeneous model collaboration, using smaller draft models for speculative decoding~\cite{leviathan2023speculative} and mixture-of-experts (MoE) routing~\cite{switch2022,mixtral2024}, can direct simple queries to lightweight models and complex queries to large ones.
The expected result is a reduction in end-to-end latency at equivalent output quality, with a secondary benefit of increased KV cache hit rates $H$.

\paragraph{Agent security architecture (L4).}
As agents acquire increasingly powerful external capabilities, including reading and writing files, executing commands, and accessing the network, the security perimeter becomes dangerously blurred.
We propose embedding sandboxing, auditing, and capability-based security directly into the agent runtime~\cite{debenedetti2025camel,ironclaw2025,agenticsecurity2025}, ensuring that every side-effecting operation undergoes explicit approval and produces an auditable event record.
The target is 100\% audit completeness: every side-effecting operation across all ICA layers must be traceable, establishing the security infrastructure required to enforce Axioms~5 and~6 in practice.

\paragraph{Weakest-link evaluation framework (full stack).}
The Barrel Principle introduced in Section~\ref{sec:agentevolution} posits that system intelligence is bounded by its weakest capability dimension, yet current evaluation practice overwhelmingly relies on aggregate or average metrics.
We propose an evaluation suite built around worst-at-$k$ scoring and multi-dimensional capability profiling (covering reasoning, safety, tool use, multilingual understanding, and physical-world interaction) that serves as a compliance test for the intelligent ISA.
The expected outcome is a compliance-grade evaluation standard for intelligent compute cores, transforming substrate independence from an architectural principle into an enforceable contractual property.

\subsection{Long-Term Research Topics}

The long-term items target fundamental architectural innovations, with the goal of constructing a truly native model-based computing architecture rather than incrementally improving today's systems.

\paragraph{Intelligent ISA and programmable skill layer (L1--L6).}
Prompts, tool descriptions, and skill packages collectively form the ``soft instructions'' of model-native computing, yet they currently lack formal semantics and stability guarantees.
We envision developing task-specific domain-specific languages (DSLs), controlled output formats, and reusable skill packages that together constitute a stable interface layer, the model-native analog of a classical Instruction Set Architecture.
As argued in Section~\ref{sec:agentevolution}, the stability of this interface layer is a prerequisite for substrate independence: whether the underlying hardware is a silicon GPU, a neuromorphic chip, or a biological compute substrate, any implementation satisfying the intelligent ISA contract should be capable of running the same agent framework.
The expected outcome is significantly improved agent generalization, enabling the same skill package to migrate seamlessly across different models and runtimes, thereby operationalizing Axioms~2 (Layered Abstraction) and~3 (Probabilistic Execution).

\paragraph{Stateful individual agents (L3--L5).}
Most agents today are stateless: each session starts from scratch, unable to accumulate or reuse experience.
We propose equipping agents with long-term memory and experience abstraction mechanisms, enabling them to learn strategies from historical executions, distill recurring patterns, and proactively optimize their own behavior~\cite{amem2025,agentmemorysurvey2025}.
The expected outcome is a measurable increase in \textit{experience utilization}, defined as the fraction of current task performance attributable to distilled prior experience, which opens new validation dimensions for both the Context Budget Law (Law~II) and the Agent Speedup Law (Law~III).

\paragraph{Distributed model-native compute fabric (L1--L6).}
When multiple agent clusters collaborate across nodes, distributed systems problems such as memory replication, state consistency, and load balancing remain largely unsolved in the model-native computing context.
Drawing on distributed database techniques for memory replication and consistency management~\cite{spanner2012,raft2014}, we propose a unified resource scheduling and state management substrate for multi-node agent clusters.
The target is for cluster utilization to scale linearly (or at least near-linearly) with node count, empirically validating the Semantic Locality Law (Law~I) and the Agent Speedup Law (Law~III) in distributed settings.

\paragraph{Governance-layer primitives (L5).}
Current governance mechanisms in agent systems (permissions, auditing, compliance) are retrofitted patches rather than first-class architectural primitives.
We propose elevating the dual-plane separation (probabilistic execution plane and deterministic control plane), audit logging, and permission declarations to core architectural primitives~\cite{arbiteros2025,ironclaw2025}, ensuring that every ICA layer has a corresponding governance mechanism.
The target is 100\% governance coverage: every side-effecting operation at every layer falls under governance constraints, fully realizing Axioms~3,~5, and~6 across the entire stack.

\paragraph{Substrate-independent intelligent compute stack (L1--L2).}
Section~\ref{sec:agentevolution} discussed the rapid development of non-silicon substrates, including neuromorphic processors, biological computing, and compute-in-memory architectures, yet current ICA interface contracts implicitly assume a GPU-centric execution model.
We propose investigating how L1--L2 interface contracts can be generalized to accommodate diverse substrates: event-driven spiking computation, analog in-memory computation, and the asynchronous response patterns of biological neural networks.
A baseline intelligence score threshold would serve as the core admission criterion.
The expected outcome is a cross-substrate interface contract specification that enables future non-silicon compute cores to integrate seamlessly with existing agent frameworks.

\paragraph{Physical-world interface drivers (L4--L5).}
Section~\ref{sec:agentevolution} identified the extension from purely digital interaction to the physical world, spanning robotics, autonomous driving, agriculture, and space navigation, as the next inflection point for intelligent operating systems, yet current L4 tool protocols (MCP, A2A) were designed for the digital domain.
We propose extending ICA's L4--L5 interfaces to support sensors, actuators, simulation environments, and telemetry data streams, following the same abstraction patterns as MCP tool servers.
The expected outcome is a tool driver specification for the physical world such that ``adding a robotic arm driver'' follows the same workflow as ``adding an MCP tool server.''

\subsection{A Note on Decoupling from Model Advances}

It is worth emphasizing that these research topics do not require stronger foundation models as a prerequisite.
Many of the problems, such as context compilation, Tool ABI standardization, and agent security architecture, are fundamentally systems engineering challenges that are orthogonal to improvements in base model capability.
This decoupling mirrors the history of classical computing: architectural breakthroughs have come not only from faster transistors but equally from more mature interfaces and runtimes.
The stability of the ISA enabled a thriving software ecosystem; the invention of virtual memory freed programs from physical memory constraints; the POSIX standard allowed applications to migrate across platforms.
Model-native computing systems stand in need of precisely the same kind of systems-level breakthroughs.

From an industry perspective, the analysis in Section~\ref{sec:sociotechnical} suggests that these systems-level standardization and decoupling efforts are co-evolving with structural changes in the industry itself.
Protocol standardization (MCP, A2A, Tool ABI) reduces inter-layer coupling, creating the technical conditions for the fabless AI industry structure predicted in Section~\ref{sec:sociotechnical}; conversely, growing industry specialization accelerates independent optimization at each layer.
This co-evolution of technology and industry structure recapitulates the central narrative of the semiconductor and personal computer industries over the past four decades.
